% Part: turing-machines
% Chapter: undecidability
% Section: introduction

\documentclass[../../../include/open-logic-section]{subfiles}

\begin{document}

\olfileid{tur}{und}{int}
\olsection{పరిచయం}

ప్రతి ప్రమేయం, చివరకు ప్రతి అంకగణిత ప్రమేయం కూడా గణనీయం కాలేదనేది
స్పష్టమే అనిపించవచ్చు. ప్రమేయాలు మరీ ఎక్కువ; వాటి ప్రవర్తన మరీ సంక్లిష్టం.
ఉదాహరణకు, రేడియోధార్మిక కణాల క్షయం నుండి నిర్వచించిన ప్రమేయాలు, లేదా ఇతర
అవ్యవస్థిత, యాదృచ్ఛిక ప్రవర్తన నుండి నిర్వచించినవి. ఒక నిర్దిష్ట క్షణం నుండి
ఒక-సెకను కాలాంతరాలను లెక్కించడం మొదలుపెట్టి, ఆ ప్రారంభ క్షణం తరువాతి
$n$వ ఒక-సెకను కాలాంతరంలో విశ్వంలోని క్షయమయ్యే కణాల సంఖ్యగా $f(n)$ను
నిర్వచించామనుకుందాం. మనం ఎప్పటికీ గణించగలమని ఆశించలేని ప్రమేయానికి ఇది
ఒక అభ్యర్థిలా కనిపిస్తుంది.

అయితే అలాంటి ప్రమేయాలను ఎలా గణించాలో ఊహించలేకపోవడం ఒక విషయం; అవి గణనీయం
కావని నిజంగా నిరూపించడం పూర్తిగా వేరొక విషయం. వాస్తవానికి, ఆశలేనంత సంక్లిష్టంగా
కనిపించే ప్రమేయాలు కూడా అమూర్తమైన అర్థంలో గణనీయాలు కావచ్చు. ఉదాహరణకు, కాలంలో
విశ్వం పరిమితమైనదనుకుందాం---కొన్ని విశ్వోద్భవ సిద్ధాంతాలు ఊహించినట్లు, సుదూర
భవిష్యత్తులో ఏదో ఒక రోజు విశ్వం ఒకే బిందువులోకి సంకోచిస్తుంది. అప్పుడు ఆ
ప్రారంభ క్షణం నుండి $f(n)$ నిర్వచితమయ్యే సెకన్ల సంఖ్య పరిమితమే (అది ఊహించలేనంత
పెద్దదైనా). పరిమిత సంఖ్యలో ఇన్‌పుట్‌లకు మాత్రమే నిర్వచితమైన ప్రతి ప్రమేయం
గణనీయం: దాని అవుట్‌పుట్‌లను ఒక పెద్ద పట్టికలో జాబితా చేయవచ్చు, లేదా ఒక అతి
పెద్ద ట్యూరింగ్ యంత్ర స్థితి-సంక్రమణ రేఖాచిత్రంలో సంకేతీకరించవచ్చు.

విలువలు అవును/కాదు ప్రశ్నలకు సమాధానాలిచ్చే ప్రమేయాల ప్రత్యేక సందర్భాలపై
మనకు తరచుగా ఆసక్తి ఉంటుంది. ఉదాహరణకు, ``$n$ ప్రధాన సంఖ్యా?'' అనే ప్రశ్నకు
కింది ప్రమేయం అనుబంధమై ఉంటుంది:
\[
\fn{isprime}(n) = \begin{cases}
  1 & \text{$n$ ప్రధాన సంఖ్య అయితే}\\
  0 & \text{లేకపోతే.}
  \end{cases}
\]
దానికి అనుబంధమైన $1/0$-విలువల ప్రమేయం ప్రభావకంగా గణనీయం అయితే, అవును/కాదు
ప్రశ్నను \emph{ప్రభావకంగా నిర్ణయించవచ్చు} అంటాం.

ప్రభావకంగా గణించలేని ప్రమేయాలు లేదా ప్రభావకంగా నిర్ణయించలేని సమస్యలు ఉన్నాయని
గణితపరంగా నిరూపించాలంటే, ఒక నిర్దిష్ట గణనా నమూనాను స్థిరపరచి, అది గణించలేని
ప్రమేయాలు లేదా నిర్ణయించలేని సమస్యలు ఉన్నాయని చూపడం తప్పనిసరి. ఉదాహరణకు,
ప్రతి ప్రమేయాన్నీ ట్యూరింగ్ యంత్రాలు గణించలేవనీ, ప్రతి సమస్యనూ ట్యూరింగ్
యంత్రాలు నిర్ణయించలేవనీ చూపవచ్చు. తరువాత చర్చ్--ట్యూరింగ్ సిద్ధాంతప్రతిపాదనను
ఆశ్రయించి, ప్రతి ప్రమేయాన్నీ గణించడానికి ట్యూరింగ్ యంత్రాల శక్తి చాలదని
మాత్రమే కాక, ఏ ప్రభావక ప్రక్రియ శక్తీ చాలదని నిర్ధారించవచ్చు.\sourcecorrection{OLTETURUND-001}{ఈ పేరాలోని ఆంగ్ల మూలం ``problems that cannot effectively decided'' అని సహాయక క్రియను వదిలేసింది. సమాంతర ప్రమేయ వాక్యానికీ తరువాతి ``problems it cannot decide'' వాక్యానికీ అనుగుణంగా, ప్రభావకంగా నిర్ణయించలేని సమస్యలు అనే ఉద్దేశిత అర్థాన్ని ఇచ్చాం.}

అలాంటి నిషేధాత్మక ఫలితాలను నిరూపించడంలో కీలకం, ట్యూరింగ్ యంత్రాలకే సంఖ్యలను
కేటాయించగలమనే వాస్తవం. దీనికి అత్యంత సులభమైన పద్ధతి వాటిని లెక్కించడం:
ట్యూరింగ్ యంత్రాలను, వాటి ప్రోగ్రాములను వ్రాయడానికి ఒక నిర్దిష్ట పద్ధతిని
స్థిరపరచి, తరువాత వాటిని క్రమబద్ధంగా జాబితా చేయవచ్చు. ఇది చేయవచ్చని గ్రహించిన
తరువాత, ట్యూరింగ్-గణనీయం కాని ప్రమేయాల ఉనికి సాధారణ కార్డినాలిటీ పరిశీలనల
నుండి వస్తుంది: $\Nat$ నుండి~$\Nat$కు వెళ్లే ప్రమేయాల సమితి (నిజానికి,
$\Nat$ నుండి $\{0, 1\}$కు వెళ్లేవాటి సమితి కూడా) \tetoken{లెక్కించలేని}{!!{nonenumerable}}; కానీ
అన్ని ట్యూరింగ్ యంత్రాలనూ లెక్కించగలం కనుక, ట్యూరింగ్-గణనీయ ప్రమేయాల సమితి
కేవలం \tetoken{అనంతంగా లెక్కించదగిన}{!!{denumerable}}.

గణనీయం కానివనీ, అనిర్ణయనీయాలనీ వరుసగా నిరూపించగల \emph{నిర్దిష్ట}
ప్రమేయాలు, సమస్యలను కూడా నిర్వచించవచ్చు. అలాంటి ఒక సమస్యే \emph{ఆగే సమస్య}.
ట్యూరింగ్ యంత్రాల నిర్దేశాలను జాబితా చేసి వాటిని పరిమితంగా వర్ణించవచ్చు.
ట్యూరింగ్ యంత్రానికి అలాంటి వర్ణనను, అంటే ట్యూరింగ్ యంత్ర ప్రోగ్రామును, మరొక
ట్యూరింగ్ యంత్రానికి ఇన్‌పుట్‌గా తప్పకుండా ఉపయోగించవచ్చు. కాబట్టి ఇతర ట్యూరింగ్
యంత్రాల గురించిన ప్రశ్నలను నిర్ణయించే ట్యూరింగ్ యంత్రాలను పరిశీలించవచ్చు.
ప్రత్యేకించి ఆసక్తికరమైన ఒక ప్రశ్న: ``ఇచ్చిన ట్యూరింగ్ యంత్రాన్ని ఇన్‌పుట్~$n$పై
ప్రారంభిస్తే అది చివరకు ఆగుతుందా?'' ఈ ప్రశ్నను నిర్ణయించగల ట్యూరింగ్ యంత్రం
ఉంటే బాగుంటుంది: ట్యూరింగ్ యంత్రాలు అనంత చక్రాల్లో చిక్కుకోకుండా చూసే నాణ్యతా
నియంత్రణ ట్యూరింగ్ యంత్రంగా దాన్ని ఊహించండి. ట్యూరింగ్ నిరూపించిన ఆసక్తికరమైన
వాస్తవం ఏమిటంటే, అలాంటి ట్యూరింగ్ యంత్రం ఉండదు. ట్యూరింగ్ యంత్రం~$M$ యొక్క
వర్ణన, ఒక సంఖ్య~$n$తో కూడిన ఇన్‌పుట్‌పై ప్రారంభించినప్పుడు, $M$ను ఇన్‌పుట్
$n$పై ప్రారంభిస్తే అది ఆగేదా లేదా అనుసరించి ఎల్లప్పుడూ $1$ లేదా $0$ అవుట్‌పుట్‌తో
ఆగే ఒకే ట్యూరింగ్ యంత్రం ఉండదు.\sourcecorrection{OLTETURUND-002}{ఆంగ్ల మూలంలోని చివరి వాక్యం ``whether $M$ machine would have halted'' అని $M$ తరువాత అనవసరంగా ``machine''ను మళ్లీ చేర్చింది. ముందే నిర్వచించిన యంత్రం~$M$ గురించే ప్రశ్న అని స్పష్టం చేశాం.}

నిర్దిష్ట అనిర్ణయనీయ సమస్యలకు ఉదాహరణలు లభించిన తరువాత, ఇతర సమస్యలు కూడా
అనిర్ణయనీయాలని చూపడానికి వాటిని ఉపయోగించవచ్చు. ఉదాహరణకు, ``మొదటిస్థాయి
\tetoken{సూత్రం}{!!{formula}}~$!A$ చెల్లుబాటవుతుందా?'' అనేది ఒక ప్రసిద్ధ అనిర్ణయనీయ సమస్య.
మొదటిస్థాయి \tetoken{సూత్రం}{!!{formula}}~$!A$ను ఇన్‌పుట్‌గా ఇచ్చినప్పుడు, $!A$ చెల్లుబాటవుతుందా
లేదా అనుసరించి $1$ లేదా $0$ అవుట్‌పుట్‌తో తప్పక ఆగే ట్యూరింగ్ యంత్రం లేదు.
చారిత్రకంగా, ఈ సమస్యను ప్రభావకంగా పరిష్కరించే ప్రక్రియను కనుగొనే ప్రశ్నను
కేవలం ``ఆ'' నిర్ణయ సమస్య అని పిలిచేవారు; అందువల్ల నిర్ణయ సమస్య పరిష్కరించలేనిదని
అంటాం. ట్యూరింగ్, చర్చ్ ఈ ఫలితాన్ని దాదాపు ఒకే సమయంలో స్వతంత్రంగా నిరూపించారు;
కాబట్టి దీన్ని చర్చ్--ట్యూరింగ్ సిద్ధాంతం అని కూడా అంటారు.

\end{document}
